\section{The noncritical rational component at conductor thirty-two}
\label{sec:conductor32-noncritical-strict}

\paragraph{Scope and original coefficient.}
Let $K=\mathbb Q(\omega)$, let $\eta_K=\chi_{-3}$, and put
\[
 L=\mathbb Q(\zeta_{96}),\qquad F=L^+,\qquad [F:\mathbb Q]=16.
 \tag{N32.1}
\]
The exact class-quotient certificate in AC-G180 proves
$\operatorname{Cl}(F)=1$.  We use that proved statement on this
same field.  No unit-basis assertion or new class computation is needed.

Fix any finite extension $E/\mathbb Q_3$ and its original ring of
integers $O$, containing the values of an even Dirichlet character
$\delta$ of conductor dividing $32$.  Write $D=E/O$.
For every positive integer $e$ define the two rational finite characters
and their full coefficient actions by
\[
 \begin{aligned}
 \psi_{\delta,e}^{\rm nc}&=\delta\eta_K^e,
 &a_{\delta,e}^{\rm nc}&=\psi_{\delta,e}^{\rm nc}c_3^{-e},\\
 \psi_{\delta,e}^{\rm cr}&=\delta\eta_K^{e+1},
 &a_{\delta,e}^{\rm cr}&=\psi_{\delta,e}^{\rm cr}c_3^{-e}.
 \end{aligned}\tag{N32.2}
\]
The actions of complex conjugation on the two full coefficients are
$+1$ and $-1$, respectively.  In particular the finite-character
parities are $(-1)^e$ and $(-1)^{e+1}$.
The word noncritical in this section means precisely this first
parity; it does not import a special-value bound.

For a number field $M$, define $\mathcal S_3(M,A)$ to be the kernel
of global continuous $H^1$ restricted to full $H^1(M_v,A)$ for
every $v\mid3$, and to $H^1(I_v,A)$ at every other finite place.
For $M=\mathbb Q$ also include the real-place $H^1$ restriction.
The latter group is zero for any $3$-primary coefficient, since
the order of complex conjugation is invertible on that coefficient.
One may equivalently use a finite localization set containing all
coefficient-ramification primes: at other primes the strict cocycle
vanishes on inertia, so it factors through that localization quotient.

\begin{theorem}[AC-G193: noncritical conductor-thirty-two strict zero]
\label{AC-G193}
For the original $O$, every such $\delta$, and every integer $e\geq1$,
\[
       \mathcal S_3\bigl(\mathbb Q,D(a_{\delta,e}^{\rm nc})\bigr)=0.
 \tag{N32.3}
\]
There is an exact decomposition on this same original coefficient:
\[
 \begin{split}
 \mathcal S_3\bigl(K,D(\delta c_3^{-e})\bigr)
 &\simeq
 \mathcal S_3\bigl(\mathbb Q,D(a_{\delta,e}^{\rm nc})\bigr)
 \oplus
 \mathcal S_3\bigl(\mathbb Q,D(a_{\delta,e}^{\rm cr})\bigr)\\
 &\simeq
 \mathcal S_3\bigl(\mathbb Q,D(a_{\delta,e}^{\rm cr})\bigr).
 \end{split}\tag{N32.4}
\]
This holds for arbitrary ramification of $O/\mathbb Z_3$.
It is not an assertion that the group on the left is zero.
In particular it applies to the two even primitive conductor-$32$
labels whose value $\beta=\overline{\delta(5)}$ in the specified
original residue field satisfies $\beta^2-\beta-1=0$.
The complementary pair with $\beta^2+\beta-1=0$ was already covered
by the stronger original strict-zero theorem AC-G180.
\end{theorem}

\begin{proof}
\proofstep{1. The unique prime of the real field.}
The powers of $3$ modulo $32$ are
\[
            1,3,9,27,17,19,25,11,1.
 \tag{N32.5}
\]
Thus $3$ has order $8$, and $-1$ does not belong to its cyclic
subgroup.  The field $\mathbb Q(\zeta_{32})$ is unramified at $3$;
it has two primes there, each of residue degree $8$.
Adjoining $\zeta_3$ gives a totally ramified quadratic local
extension.  It is disjoint from the unramified degree-eight local
extension, and hence $L$ has two primes above $3$, both with
ramification index $2$ and residue degree $8$.

Under the product identification
\[
 \operatorname{Gal}(L/\mathbb Q)
  \simeq (\mathbb Z/3\mathbb Z)^\times
             \times(\mathbb Z/32\mathbb Z)^\times,
 \tag{N32.6}
\]
a decomposition subgroup at $3$ is the product of the full first
factor and $\langle3\bmod32\rangle$ in the second factor.
Complex conjugation is $(-1,-1)$, so it is not in that subgroup.
It exchanges the two primes.  Their common restriction to $F=L^+$
is therefore the unique prime $v$ of $F$ above $3$.
The extension $L/F$ splits at $v$, and
$[F_v:\mathbb Q_3]=16$, with $e_v=2$ and $f_v=8$.
These assertions use the ordinary unramified Frobenius description
of a prime in a cyclotomic extension of conductor prime to it,
and the Eisenstein polynomial of $\zeta_3-1$ at $3$.

\proofstep{2. One-prime local generation over a real field at an odd prime.}
We give the needed extension of the proof of AC-G182.
Let $M$ be any number field, let $p$ be odd, assume that $M$ has
one place $v$ above $p$, and assume $\operatorname{Cl}(M)[p]=0$.
Let $\Gamma$ be the Galois group of the maximal extension
unramified at every finite place other than $v$, and let
$P=\Gamma^{(p)}$ be its maximal pro-$p$ quotient.
Then the image of $G_{M_v}$ in $P$ is all of $P$.

Indeed that image is compact and thus closed.  If proper, it is
proper in some finite $p$-group quotient $Q$ of $P$.
A proper subgroup of a finite $p$-group is contained in a normal
subgroup of index $p$: for a maximal subgroup this follows by
induction from the nontrivial center, and every proper subgroup
is contained in a maximal one.  The resulting quotient $C_p$
gives a nontrivial cyclic extension of $M$, unramified away $v$
and split at $v$.  It is unramified at every finite place.
It is also split at every real place.  A real decomposition group
has order at most two, so its image in the odd-order group $C_p$
is trivial.  Ordinary Hilbert class-field reciprocity now gives
a quotient of $\operatorname{Cl}(M)$ of order $p$, a contradiction.
This explicitly handles the real-place clause omitted from the
totally imaginary formulation of AC-G182.

\proofstep{3. Cohomology with the full local witness.}
Let $O'$ be any finite extension of $\mathbb Z_p$ with fraction
field $E'$, and let $a:\Gamma\to1+\mathfrak m_{O'}$ be continuous.
Put $A=(E'/O')(a)$.  A continuous cocycle $z:\Gamma\to A$
has finite image.  If its local cohomology class is zero, choose
one local witness $d\in A$ with $z(g)=(a(g)-1)d$ on $G_{M_v}$.
Choose $m$ so that both this witness and the entire cocycle image
belong to $A[\pi^m]$, for an original uniformizer $\pi$.
The action on that finite module has finite $p$-group image $U_m$,
because it comes from $1+\mathfrak m_{O'}$.
The cocycle homomorphism
\[
       g\longmapsto (z(g),a(g)|_{A[\pi^m]})
                \quad\hbox{in }A[\pi^m]\rtimes U_m
 \tag{N32.7}
\]
therefore factors through $P$.
Step 2 makes its local and global images equal.
For each global $g$, choose a local element with the same image
in (N32.7).  Since the chosen witness is in the same finite
module, its local coboundary formula gives
$z(g)=(a(g)-1)d$ globally.  Thus restriction on $H^1$ is injective.

The assertion applies to a full strict class on $G_M$ when $a$
is unramified outside $v$.
At any other finite place inertia acts trivially on $A$, so a
zero inertia cohomology class means the cocycle itself is zero
on that inertia group.  The finite homomorphism (N32.7) is then
trivial on their closed normal closure and factors through
$\Gamma$.  The same witness still works.  Consequently
\[
 \mathcal S_p\bigl(M,(E'/O')(a)\bigr)=0.
 \tag{N32.8}
\]
No invariant-divisibility argument or class-group inverse-limit
control is used here.

\proofstep{4. The precise noncritical action over $F$.}
Apply Steps 2--3 with $M=F$, $p=3$, and the unchanged original $O$.
Step 1 supplies its unique prime and AC-G180 supplies its class group.
The image of $G_F$ in $\operatorname{Gal}(L/\mathbb Q)$ is
$\{1,c\}$.  Because $\delta$ is even and factors through $L$,
its restriction to $G_F$ is trivial.
The residual cyclotomic character at $3$ is exactly
$\overline{c_3}=\chi_{-3}$ on $G_{\mathbb Q}$; this follows from
the action on $\mu_3$.  Thus
\[
 a_{\delta,e}^{\rm nc}|_{G_F}
             = (\eta_K/c_3)^e|_{G_F}
             \quad\hbox{takes values in }1+\mathfrak m_O.
 \tag{N32.9}
\]
Both factors on the right are unramified outside $3$.
The restricted coefficient therefore satisfies (N32.8), giving
$\mathcal S_3(F,D(a_{\delta,e}^{\rm nc}))=0$.

A rational strict class restricts into this strict group:
full local zero at $3$ restricts to full local zero at $v$,
and zero on any other inertia group restricts to zero on every
inertia subgroup above it.  Restriction followed by corestriction
on $H^1$ is multiplication by $[F:\mathbb Q]=16$.
This is an automorphism of a $3$-primary group, hence restriction
is injective.  This proves (N32.3) on the original coefficient.

\proofstep{5. The exact index-two decomposition of full strict groups.}
The two characters in (N32.2) restrict to the same coefficient
$\delta c_3^{-e}$ on $G_K$, since $\eta_K|_{G_K}=1$.
For their original integral lattices the induced representation
from $G_K$ splits as
\[
 \operatorname{Ind}_{G_K}^{G_{\mathbb Q}}O(\delta c_3^{-e})
       \simeq O(a_{\delta,e}^{\rm nc})
                      \oplus O(a_{\delta,e}^{\rm cr}).
 \tag{N32.10}
\]
In a two-coset basis its change-of-basis matrix has determinant
a unit times $2$, which is invertible in this $O$.
Equivalently the two idempotents $(1\pm\tau)/2$ act integrally
on the regular $C_2$ factor.  Tensor with $E/O$ and apply Shapiro.

At $3$, local Shapiro identifies the full local group with the
group at the unique place of $K$ above $3$.
At every other rational finite prime the extension $K/\mathbb Q$
is unramified.  The local Mackey decomposition, restricted to
inertia, therefore identifies zero inertia restrictions with
zero on the actual inertia groups at all places of $K$ above it.
For a split prime both places occur; for an inert prime their
conjugate inertia components give the same zero condition.
Thus this is an equality of the full strict kernels, not just
an isomorphism of unrestricted global cohomology.
At infinity induction is from the complex place of $K$;
all real $H^1$ groups with these $3$-primary coefficients are zero.
This proves the first isomorphism in (N32.4), and Step 4 proves
the second.  There is no finite index or coefficient-rank error.
\end{proof}

\paragraph{Boundary of the result.}
This theorem supplies only the vanishing of the noncritical
rational component and the exact comparison (N32.4).
Any upper bound for its critical component is a separate input.
In particular AC-G191 and an actual-window composition must be
independently checked before this lemma can remove two terms
from a main arithmetic remainder or change its sufficient gate.
No such window assertion, class-tower control theorem, or mixed
theta/Selmer realization is claimed in this section.

\paragraph{Source scope.}
The field and class-group input is the exact AC-G180 real-$96$
certificate, including its class-quotient-only interpretation.
The finite pro-$p$ and cocycle argument is repeated in full from
the method of AC-G182, with the real-place extension proved above.
Its retained deep input is ordinary Hilbert class-field reciprocity
(Sharifi, \emph{Iwasawa Theory}, Theorem 1.1.11 and preceding
definition); no new proof of that theorem is asserted.
Restriction/corestriction, Shapiro and local Mackey are the same
standard continuous-cohomology tools used in the existing exact
coefficient and local comparisons.  Their concrete index-two
matrix and all local conditions are specified here.
